Le dihydrogène brûle dans le dioxygène selon l'équation balancée :
\[ 2\,\text{H}_2 + \text{O}_2 \rightarrow 2\,\text{H}_2\text{O} \]
On fait réagir complètement \qty{\v{mH}}{g} de dihydrogène. Masses molaires : $M(\text{H}_2) = \qty{2.016}{g/mol}$, $M(\text{O}_2) = \qty{32.00}{g/mol}$, $M(\text{H}_2\text{O}) = \qty{18.02}{g/mol}$.

\begin{enumerate}[label=\alph*)]
    \item Combien de moles de dihydrogène réagissent ?
    \item Quelle masse de dioxygène est nécessaire ?
    \item Quelle masse d'eau est produite ? Vérifiez que la loi de conservation de la masse est respectée.
\end{enumerate}
